% Expository adaptation, October 7, 2026. Original mathematics: OpenAI, September 23, 2026.
% Prose revised and explanations expanded; see README.md for the source and scope.
\section{From the estimator to deterministic simulation}
\label{sec:conclusion}

We now put the three parts together. The exact hierarchy supplies a reward close to acceptance probability. The estimator approximates that reward for most environments. The controller evaluates every environment in logarithmic space and always halts. These statements have different quantifiers, and the final reduction needs each in its stated form.

There are three errors to allocate: the residual transition contribution, the estimator's statistical error at the designated start, and truncation of its numerical digits. We choose the parameters so that their sum is below the requested accuracy. Enumerating the conditional environments and taking a median then produces a deterministic approximation, without a promise that the acceptance probability lies outside a decision gap.

All transducers below have read-only input and a one-way write-only
output tape. Output contents are never available as work memory.
For a fixed randomized machine $\mathcal M$, write
$p_{\mathcal M}(x)$ for its acceptance probability on input $x$.

\begin{theorem}[Fixed inverse-polynomial accuracy]
\label{thm:fixed-precision}
Fix a randomized machine $\mathcal M$ with polynomial worst-case
running time and $O(\log(n+2))$ work space, and fix an integer $d\ge1$.
There is a uniform deterministic transducer that, on every input $x$
of length $n$, outputs a dyadic rational $z_x\in[0,1]$ with
\[
 |z_x-p_{\mathcal M}(x)|\le(n+2)^{-d}.
\]
Its work space and output length are $O(\log(n+2))$, and its running
time is polynomial in $n+2$. The constants may depend on
$\mathcal M$ and $d$, but not on the input.
\end{theorem}

The fixed exponent in this theorem is enough to decide bounded-error
computations. After proving it, we derive the class equality and its
promise version. We then allow the requested accuracy to be part of the
input and construct accepting computations under an inverse-polynomial
probability promise.

\subsection{Choose constants before choosing the input-dependent ranges}
\label{subsec:final-constants}

Several parts of the proof ask for a fixed constant to be sufficiently large or small. We record their order here to ensure that no later choice invalidates an earlier estimate. In particular, the averaging gap is uniform in the field size and number of free bits; the field can therefore be chosen after the required identifier and sampling ranges are known.

The construction has the following order of choices.
\begin{enumerate}
\item Fix $H,h_0,s=8$ and the fixed number of independent channels.
The finite-action gap is uniform in every later modulus and bit length.
\item Choose the correction thresholds and cutoff, then the iteration
count $k_*$ from the factorial bound and its resulting uniform
magnitude constants. These estimates hold for every column multiplier
in $[0,1]$.
\item Choose $D_0$ for the density reduction, the routing thresholds,
and the finite ladders of entry and gate thresholds. The row-gate
rank $K$ can then be chosen from the row-mass bounds alone.
\item Choose the detector rank $J$ large enough for the incidence
bound on extra copies. Choose $R_*$ large enough for overload
detection, and finally the incoming rank $K_c$ large enough for
correction normalization. The earlier row-gate choice does not depend
on $K_c$.
\item Fix the bounded-output constant, the scale cutoff and precision
coefficient, and the prefix mean exponent. Choose the buffer
floor and capacity exponents and its noise coefficient so that all
buffer sums converge. Choose task separation and walk length
coefficients sufficiently large for the small constants in the
clean estimation induction.
\item Fix the crude range and propagation constants, then the block
gain coefficients $D_2,D_1$ and a sufficiently small rational block
width coefficient $\epsilon$. These choices close the coupled
fingerprint induction of Theorem~\ref{thm:fp-estimators}.
\item Choose the fixed operation-stratum and allowance coefficients
and the guard constants in Theorem~\ref{thm:ctrl-realization}.
Their finite enlargement does not alter the analytical tables or
the distribution of the environment.
\item Choose the linear top-budget coefficient, the field-size
exponent, and the minimum padding of $B$.
\end{enumerate}
Within each step only finitely many strict inequalities between
fixed constants are needed. The proofs in the corresponding sections
give these inequalities. In particular, the row rank is fixed before
the incoming rank, the block constants follow the crude error
constants, and the field is chosen after the reachable resource
ranges. There is no dependence of an earlier sampler gap on that
last choice.

The coefficient $\epsilon$ in this order is the fixed fingerprint
block-width coefficient. Once it has been chosen, increasing the fixed
linear top budget increases the number of salt blocks but leaves the
sampler's gap uniform: all earlier salt bits form one free component in
Lemma~\ref{lem:env-gap}. The storage estimate
\eqref{eq:fp-salt-space} therefore remains $O(B)$.

For clarity, one may bound the stage index by $L$, the inherited
suffix length by the stages remaining above a task, and logarithms
of reciprocal requested rates by
\begin{equation}\label{eq:final-admissible}
HL+c(M_{\rm top}-M)
\end{equation}
with a sufficiently large fixed $c$. A child at slack $d$ can extend
the variable scale range by only $O(d)$ and loses at least $d$ of
budget. Stage-rate resets are already covered by $HL$.
The finitely many local task operations use the stated widths and
fixed subranks. Thus Equation~\eqref{eq:final-admissible} is closed
under all estimator recipes. The controller guards allow these
designated executions.

There is consequently a fixed $C_P$ such that a prime
\[
2^{C_PB}\le P\le 2^{C_PB+1}
\]
is larger than all identifiers, makes the radix invertible, and gives
$P2^{-k}\ge1$ for every rate ever requested. Bertrand's postulate
provides such a prime; testing consecutive integers by trial division
finds one using $O(B)$ space. The search range and running time are
$2^{O(B)}$. Field arithmetic and fixed-dimensional rank tests also
use $O(B)$ space.

\subsection{Pass from average error to the specified start}

The estimator theorem averages over source vertices. We need the value at one particular vertex, so we cannot use that theorem without paying for the loss of averaging over vertices. The next elementary calculation keeps one nonnegative summand and makes the cost explicit. Conditioning on the start being retained also identifies the exact family of environments that the deterministic algorithm will enumerate.

\begin{lemma}[Conditioning on one start]\label{lem:final-majority}
Let $W\in[0,1]^V$, and suppose an estimator satisfies
\[
\frac1N\mathbb E\sum_{v\in V}\frac{X_f(v)}f
 |\bar W(v)-W(v)|^8\le C_b^8 2^{-8p}.
\]
For any specified $z\in V$ and every $\eta>0$,
\begin{align}
\mathbb E\bigl[|\bar W(z)-W(z)|^8\mid X_f(z)=1\bigr]
 &\le N\frac f{\pi_f}C_b^8 2^{-8p},
 \label{eq:final-conditional}\\
\Pr\bigl[|\bar W(z)-W(z)|>\eta\mid X_f(z)=1\bigr]
 &\le N\frac f{\pi_f}C_b^8 2^{-8p}\eta^{-8}.
 \label{eq:final-tail}
\end{align}
\end{lemma}

\begin{proof}
Keep only the nonnegative summand at $z$. Since
$\Pr[X_f(z)=1]=\pi_f>0$,
\[
 \frac{\pi_f}{f}
 \mathbb E\bigl[|\bar W(z)-W(z)|^8\mid X_f(z)=1\bigr]
 =\mathbb E\!\left[\frac{X_f(z)}f|\bar W(z)-W(z)|^8\right]
 \le N C_b^8 2^{-8p}.
\]
Dividing by $\pi_f/f$ proves
Equation~\eqref{eq:final-conditional}. In particular, choosing one
specified coordinate costs the factor $N$ from the integrated norm.
Markov's inequality gives Equation~\eqref{eq:final-tail}.
\end{proof}

\begin{proof}[Proof of Theorem~\ref{thm:fixed-precision}]
We choose the hierarchy depth and statistical budget first, then convert each estimator value to a finite dyadic number. After proving that most such numbers are accurate, we explain how to find their median and count the total memory and time.

Apply Lemma~\ref{lem:model-reduction} to $\mathcal M$, with no promise
on its acceptance probability. Set
\[
 r=d\lceil\log_2(n+2)\rceil,\qquad \Delta=2^{-r},\qquad
 L=\left\lceil\frac{C_0B+r+3}{H-1}\right\rceil,
\]
where $|V_0|\le2^{C_0B}$ is an explicit configuration bound.
Since $d$ is fixed, $r,L=O(B)$. Let $z$ be the all-primary copy of
the start at stage $L$. Proposition~\ref{prop:hier-reward-target}
gives
\begin{equation}\label{eq:final-hierarchy-accuracy}
 0\le p_{\mathcal M}(x)-W_L(z)\le\Delta/8.
\end{equation}
This bound holds for every acceptance probability. Choose a fixed
$C_N$ large enough to bound all the resulting stage domains by
$N=2^{C_NB}$.

For the final reward task take $a=u=f=b_L$. Its rank satisfies
$i_*=O(L+1)$, and Equation~\eqref{eq:est-precision} reads
\[
 p=M_{\rm top}/K_{\rm prec}-HL-A_*i_*.
\]
By Lemma~\ref{lem:env-ranks}, fix an integer $C_{\rm tail}\ge1$ bounding
$(f/\pi_f)C_b^8$ throughout the admissible range. Choose the fixed
linear coefficient of $M_{\rm top}=O(B)$ large enough that
\begin{equation}\label{eq:final-precision-choice}
 p\ge r+3+\frac{C_NB+\lceil\log_2 C_{\rm tail}\rceil+3}{8}.
\end{equation}
This is the extra accuracy needed to pay the factor $N$ in the passage to a specified vertex. The top-budget and field-size choices in the preceding subsection allow this enlargement; Equation~\eqref{eq:final-admissible} still
bounds every descendant rate. Theorem~\ref{thm:fp-estimators}
supplies the reward moment estimate. Lemma~\ref{lem:final-majority},
with $\eta=\Delta/8$, now shows that
\begin{equation}\label{eq:final-good-environments}
 \Pr\bigl[|\bar W_L(z)-W_L(z)|>\Delta/8
       \mid X_f(z)=1\bigr]\le1/8.
\end{equation}
Thus the factor $N$ lost on selecting one coordinate has been paid
by the statistical precision. Explicitly, the bound in
Equation~\eqref{eq:final-tail} is at most
\[
 C_{\rm tail}\,2^{C_NB-8p+8r+24}\le 2^{-3},
\]
by Equation~\eqref{eq:final-precision-choice}.

We next turn the estimator values into a finite list of dyadic numbers.
For each master environment, Theorem~\ref{thm:arith-interface}
gives the digits $w_j\in\{-6,\ldots,6\}$ of its one fixed value
$\bar W_L(z)$. Structural choices do not depend on the requested
digit. Put $K=r+6$, compute
\[
 A_\sigma=\sum_{j=0}^{K}w_j2^{K-j},\qquad
 a_\sigma=\min\{2^K,\max\{0,A_\sigma\}\},\qquad
 v_\sigma=a_\sigma2^{-K}.
\]
The integer $A_\sigma$ is obtained by starting at zero and applying
$A\leftarrow2A+w_j$ in order, so it uses $O(K)$ bits.
By Equation~\eqref{eq:arith-tail}, the truncation error is at most
$6\,2^{-K}<\Delta/8$. Clipping cannot increase error relative to
$W_L(z)\in[0,1]$. Equations~\eqref{eq:final-hierarchy-accuracy}
and~\eqref{eq:final-good-environments} imply that, on more than
three quarters of the environments conditioned on $X_f(z)=1$,
\begin{equation}\label{eq:final-dyadic-accuracy}
 |v_\sigma-p_{\mathcal M}(x)|<3\Delta/8.
\end{equation}

The majority statement refers to a particular conditional distribution, so the deterministic enumeration must reproduce that distribution exactly. Enumerate the conditional law exactly. Enumerate each field-matrix
entry in $\{0,\ldots,P-1\}$ and each raw salt bit string, retaining
precisely those encodings whose positive-degree columns form a frame
and for which $X_f(z)=1$. The retained encodings are equally weighted
and give the law in Equation~\eqref{eq:final-good-environments}.
In particular, enumerate the original salt strings, not their reduced
fingerprint parameters. Different raw strings that induce the same
fingerprint remain separate encodings; their multiplicities reproduce
the original probability law. At least one encoding is retained: a frame
exists, and its independent constant column can place the designated
evaluation in the nonempty retention box. The matrices have fixed
dimensions and $\log P=O(B)$; together with
Equation~\eqref{eq:fp-salt-space}, this gives $O(B)$ bits per master
environment and at most $2^{O(B)}$ encodings to enumerate.

We find the median by counting, not by storing or sorting the estimates. Let $E\ge1$ be the number of retained encodings, found by one pass,
and let $h=\lfloor E/2\rfloor+1$. Binary-search the integer interval
$[0,2^K]$ for the least $a$ such that at least $h$ retained encodings
satisfy $a_\sigma\le a$. Each count is computed by a fresh pass,
recomputing each numerator $a_\sigma$. This finds the upper median
of the finite multiset, including all repetitions. Since more than
three quarters of its members satisfy
Equation~\eqref{eq:final-dyadic-accuracy}, so does its median: fewer than
$E/4$ values can lie below the indicated interval, and fewer than
$E/4$ can lie above it. The algorithm retains only the current count
and search registers; it never stores the multiset.
Output the integers $(a,K)$ representing $z_x=a2^{-K}$. Then
\[
 |z_x-p_{\mathcal M}(x)|<3\Delta/8
       \le(n+2)^{-d}.
\]

We account for the memory held during this computation. The master
environment, current cursor, catalyst, shared scratch, and suspended
controller records together use $O(B)$ bits by
Theorem~\ref{thm:ctrl-realization}, because $M_{\rm top},K=O(B)$.
The outer computation keeps only a fixed number of additional
$O(B)$-bit registers: its top parameters and designated start,
the environment enumeration index, $E$, $h$, the current count,
the two binary-search endpoints and candidate, the digit index,
and the signed numerator accumulator. Each count is bounded by the
number of environment encodings, so it also has $O(B)$ bits.
These outer registers remain live while a digit query runs; in
particular, both its environment-enumeration state and its numerator
accumulator are included. They are present once for the whole
computation, not once per recursive level. Every query terminates on every environment,
including the inaccurate environments, by the same realization
theorem. The output is written only after the median is computed.

To bound time, use totality together with the complete configuration count. This is why termination on inaccurate environments was needed, rather than only on the majority of accurate ones. Each digit query is a halting deterministic computation with $O(B)$
work bits and physical input heads confined between the real endmarkers.
Counting its finite control, work-tape contents, and all work and input
head positions gives at most $2^{O(B)}$ complete configurations on its
fixed input. A repeated complete configuration would
force a deterministic loop, so totality bounds the query by
$2^{O(B)}$ steps. There are
$2^{O(B)}$ environment encodings, $O(K+1)$ binary-search passes,
and $K+1$ digit queries per numerator. Their product is still
$2^{O(B)}=(n+2)^{O(1)}$. All finite parameters are generated
uniformly, and all constants are fixed for $(\mathcal M,d)$.
\end{proof}

\subsection{The class equality and promise separation}

\begin{proof}[Proof of Theorem~\ref{thm:main}]
Let $\mathcal M$ recognize a language in $\BPL$, with acceptance
probability at most $1/3$ on no inputs and at least $2/3$ on yes
inputs. Apply Theorem~\ref{thm:fixed-precision} with $d=3$.
For every input length, including $n=0$, its error is at most
$(n+2)^{-3}\le1/8$. Hence the dyadic answer is at most
$1/3+1/8=11/24$ on a no input and at least
$2/3-1/8=13/24$ on a yes input. Comparing this answer with $1/2$
decides the language uniformly in logarithmic space. This proves
$\BPL\subseteq\Lclass$ directly from probability approximation.

A deterministic machine may ignore its coins, giving
$\Lclass\subseteq\RL$. Two independent repetitions of an $\RL$
machine, accepting if either repetition accepts, have acceptance
probability zero on no inputs and at least $3/4$ on yes inputs.
They use logarithmic space and polynomial time, so
$\RL\subseteq\BPL$. The three inclusions prove
$\Lclass=\RL=\BPL$.
\end{proof}

A promise problem is a pair $(Y,N)$ of disjoint sets of inputs.
Here $\mathsf{PromiseBPL}$, denoted $\mathsf{prBPL}$ in the introduction,
requires one polynomial-time randomized
logarithmic-space machine that halts on every input, accepts each
input in $Y$ with probability at least $2/3$, and accepts each
input in $N$ with probability at most $1/3$.
The class $\mathsf{PromiseL}$ consists of pairs separated by a total
deterministic logarithmic-space decider; there is no correctness
requirement outside $Y\cup N$.

\begin{corollary}\label{cor:promise-bpl}
Under these conventions,
$\mathsf{PromiseBPL}=\mathsf{PromiseL}$.
\end{corollary}

\begin{proof}
The same $d=3$ approximation and threshold $1/2$ used in the preceding
proof accept every input in $Y$ and reject every input in $N$.
The resulting machine halts also outside the promise because its
approximation queries are total. The reverse inclusion follows by
ignoring the coins.
\end{proof}

\subsection{Handle an accuracy request supplied as input}

The fixed-accuracy theorem allows constants to depend on a fixed inverse-polynomial exponent. Theorem~\ref{thm:probability} instead needs one algorithm whose constants do not depend on the input accuracy request.

Merely asking for more digits would not achieve this: it approximates the same estimator more accurately while leaving its statistical error unchanged. We instead represent the computation as a padded input to one fixed inverse-polynomial approximator. The padding length increases the size parameter used by that approximator, while its bits can be generated on demand in the desired space.

\begin{proof}[Proof of Theorem~\ref{thm:probability}]
Fix a padded machine $\mathcal M_{\rm pad}$ as follows. On a string
beginning with $1^n0x$, where $|x|=n$, it simulates $\mathcal M(x)$
and ignores all remaining symbols. If the delimiter or the following
$n$ symbols are absent, it rejects. This is one fixed randomized
machine with polynomial worst-case running time and logarithmic work
space in its full input length. On every string with the specified
prefix its acceptance probability is $p_{\mathcal M}(x)$.
Fix once and for all the transducer of
Theorem~\ref{thm:fixed-precision} for $\mathcal M_{\rm pad}$ and $d=1$.

On input $(x,1^q)$, with $n=|x|$ and $q\ge1$, set
\[
 m=2^{q+\lceil\log_2(2n+4)\rceil+4}.
\]
Simulate that fixed transducer on the virtual length-$m$ string
\[
 y=1^n0x0^{m-(2n+1)}.
\]
An input query is answered by comparing its address with the prefix
boundaries and either returning the appropriate bit of $x$ or a
fixed bit. The simulated input-head address and all these comparisons
use $O(\log m)$ bits; the string $y$ is never stored. The simulated
transducer returns a dyadic rational $z\in[0,1]$ satisfying
\[
 |z-p_{\mathcal M}(x)|\le(m+2)^{-1}<2^{-q-4}.
\]
Its numerator and exponent have $O(\log m)$ bits and can be retained
in work space while its finite output is simulated.

Compute by exact dyadic arithmetic
\[
 a=\left\lfloor 2^{q+2}z+\tfrac12\right\rfloor.
\]
Then $0\le a\le2^{q+2}$, and rounding gives
\[
 \left|a2^{-(q+2)}-p_{\mathcal M}(x)\right|
 \le2^{-q-3}+2^{-q-4}<2^{-q}.
\]
Write the binary numerator $a$ to the output tape. It has at most
$q+3$ bits and represents the denominator specified in the theorem.
The simulator state, virtual input address, actual input counters,
captured dyadic answer, and rounding operands occupy a fixed number
of records of total size
\[
 O(\log m)=O(\log(n+2)+q).
\]
The simulation and the synthesis of each input bit take polynomial
time in $m$, so the total time is
$m^{O(1)}=(n+2)^{O(1)}2^{O(q)}$.
The padded machine and its exponent-one approximator were fixed
before reading $q$. Thus their constants are independent of the
requested accuracy. Throughout their simulation, the original size
convention is $B=\Theta(\log(m+2))$, exactly as required by the
estimator and controller theorems.
\end{proof}

\subsection{Construct an accepting computation from continuation values}

We return to fixed inverse-polynomial accuracy to construct an accepting
computation when the machine's acceptance probability has an
inverse-polynomial lower bound.
For a transducer, the conclusion concerns outputs certified by an
accepting run: its output tape must be write-only, so past emitted
symbols do not influence later transitions.

\begin{corollary}[Certified accepting computations]
\label{cor:certified-search}
Fix a randomized logarithmic-space machine $\mathcal M$ that halts
within a fixed polynomial bound $T(n)\ge1$ on every input and every
coin sequence, and fix an integer $c\ge0$. There is a uniform
deterministic logarithmic-space, polynomial-time transducer that
outputs an accepting computation of $\mathcal M$ whenever
\[
 p_{\mathcal M}(x)\ge(n+2)^{-c}.
\]
On every other input it outputs either an accepting computation or
a failure symbol. If every accepting run of $\mathcal M$ emits a
valid witness for a specified relation, the transducer can instead
output such a witness with the same guarantee.
\end{corollary}

\begin{proof}
The decision approximation can also guide a path. At each configuration, some successor has continuation probability at least the current weighted average. Approximate those successor probabilities closely enough that the loss incurred by choosing the largest approximation is small over the entire path.

Give halted configurations constant self-transitions through time
$T(n)$. A configuration and its elapsed time use $O(\log(n+2))$
bits. There is a fixed randomized logarithmic-space continuation
machine which, on a valid encoding of $(x,v,t)$, starts from
configuration $v$ and simulates the remaining $T(n)-t$ steps with
fresh coins. It checks the fixed work-space box and time bound,
rejecting malformed encodings and any attempted violation. Its
acceptance probability on the reachable configurations is the
continuation value $p_t(v)$, and it has a
polynomial worst-case time bound in its input length.

Write $\gamma=(n+2)^{-c}$. Choose a fixed exponent $d$ large enough
that Theorem~\ref{thm:fixed-precision}, applied to the continuation
machine, approximates every queried value within
$\delta\le\gamma/(8T(n))$. Such a fixed exponent exists because
$T$ is a fixed polynomial and the continuation input contains $x$.
The extra fields have logarithmic length. We supply them virtually
from saved configuration and time registers, so the approximation
and those registers together use logarithmic space.

Starting at the initial configuration, at each time estimate the
continuation value of every positive-probability successor and
choose one with maximum estimated value, breaking ties by a fixed
port order. Since $p_t(v)$ is the weighted mean of the successor
values, the chosen successor $v'$ satisfies
\[
 p_{t+1}(v')\ge\max_w p_{t+1}(w)-2\delta
                  \ge p_t(v)-2\delta.
\]
On an input satisfying the stated lower bound, after $T(n)$ choices
the terminal continuation value is at least
$\gamma-2T(n)\delta\ge3\gamma/4>0$. A terminal continuation value
is zero or one, so this computation accepts.

A low-probability input need not yield a successful path, so do not emit a purported accepting computation before checking where the deterministic choices lead. First run this construction without writing output. If its terminal
state rejects, output the failure symbol. Otherwise repeat the same
deterministic choices and stream the configurations through the original
first halting state, or the symbols emitted by the simulated transducer.
The artificial self-transitions after halting are omitted from the output.
This second pass follows the same accepting computation. Consequently no rejecting computation
is output, even outside the probability promise.

At any time the current, candidate, and best successor configurations,
the time counter, a constant number of dyadic scores, the virtual
input address, and the approximation workspace occupy a constant
number of $O(\log(n+2))$-bit records. The finite successor lists
have constant size, and there are only $T(n)$ steps in each pass,
so the total running time is polynomial. Each approximation has a
deterministic guarantee at its queried configuration; no simultaneous
accuracy event for all adaptive queries is needed.
\end{proof}

